\documentclass[11pt]{clapps}  % Praxia/Kontablo house style (clean-room)

% --- PACKAGES ---
% geometry, lmodern, microtype, xcolor and hyperref are provided by clapps.cls.
\usepackage{booktabs}
\usepackage{amsmath}
\usepackage{amssymb}
\usepackage{float}
\usepackage{url}
\usepackage{graphicx}          % \resizebox, used by clapps.cls's \fitwidth
\usepackage{tikz}
\usetikzlibrary{arrows.meta, positioning, fit, backgrounds}
\DeclareUrlCommand\pathtt{\urlstyle{tt}}
% NOTE on \pathtt: it is a \DeclareUrlCommand, so its argument is verbatim.
% Never write \_ inside it — the backslash prints literally. Use a bare _.

% Float placement: the two TikZ figures are tall. Without these, LaTeX's
% default topfraction/textfraction defer both to the end of the document
% instead of placing them near their references.
\renewcommand{\topfraction}{0.88}
\renewcommand{\bottomfraction}{0.7}
\renewcommand{\textfraction}{0.08}
\renewcommand{\floatpagefraction}{0.75}

% --- PREAMBLE ---
\title{Deterministic Auditability Invariants for Autonomous
Financial Agents: The Loss Ledger and Pre-Transaction Fiber Query}
\author{Christian Luciani}
\date{\small{Spoke preprint v0.2 --- July 2026 --- draft for arXiv cs.MA}}

\begin{document}

\maketitle

\noindent\textbf{Author.} Christian Luciani --- Chief Architect, Kontablo
Project. Independent Researcher, Cuenca, Ecuador (Praxia, the initiative's
planned entity, in incorporation). ORCID:~\href{https://orcid.org/0000-0002-6955-5384}{0000-0002-6955-5384}.
Contact:~\href{mailto:cluciani@gmail.com}{cluciani@gmail.com}.

\medskip

\noindent\textbf{Normative reference.} This spoke paper is a community-facing
result built on Kontablo, whose specification of record is the canonical
preprint, Zenodo concept DOI
\href{https://doi.org/10.5281/zenodo.20738795}{10.5281/zenodo.20738795}
(cite this DOI; it resolves to the latest version). All architectural
principles referenced by number (e.g.\ ``principle~\#5'') are defined there.
This paper does not restate that thesis; it isolates one question that belongs
to the multi-agent-systems community and answers it with a new, independently
regenerable result.

\begin{abstract}
\noindent Kontablo is a graph-based, UUID-keyed universal accounting ontology
that acts as a deterministic semantic-resolution layer between local statutory
charts of accounts and the machine-to-machine (M2M) agentic economy. This paper
asks a question specific to autonomous multi-agent finance: as AI agents begin
to execute high-frequency financial transactions over protocols such as AP2
(Agent Payments Protocol) and A2A (Agent2Agent), \emph{what deterministic
auditability invariants must the semantic-resolution layer satisfy so those
transactions are trustworthy without a human in the per-transaction critical
path?} We argue the answer cannot be procedural (human review does not scale to
M2M frequency) and must instead be \emph{architectural}. We name three
invariants and demonstrate them in a committed reference implementation:
(\textbf{I1}) \emph{ontology-as-constraint} --- the agent-facing interface
exposes only deterministic tools, and no tool can emit an account identifier
that does not exist in the graph; (\textbf{I2}) the \emph{loss ledger} --- every
piece of information the layer cannot translate with certainty becomes a typed,
countable record that escalates to a human, never a silent guess; and
(\textbf{I3}) the \emph{pre-transaction fiber query} --- an agent can retrieve
the exact preimage of any consolidated node \emph{before} acting on it, making
auditability a pre-condition of the transaction rather than only a forensic
afterthought. We instantiate the invariants over an agent-native
Model Context Protocol (MCP) surface of six deterministic tools and a
round-trip audit that, over $441$ ledger entries drawn from $75$ synthetic
entities across $68$ jurisdictions, reconstructs every local trial balance
byte-for-byte from lineage alone and records \emph{zero} silent losses. The
contribution to the agent community is a concrete shape for machine-verifiable
trust at the resolution layer that complements, rather than duplicates,
settlement-layer protocols.

\medskip
\noindent\textbf{Keywords:} multi-agent systems, agent-to-agent protocols, M2M
economy, deterministic auditability, provenance, Model Context Protocol,
AP2, financial agents, ontology-as-constraint.
\end{abstract}

\compactcontents

% =====================================================================
\section{Why does trust have to move into the architecture?}
% =====================================================================

An autonomous agent that pays another agent must, at some point, commit a
financial fact: \emph{this outflow of value corresponds to this economic event,
booked to this account, in this entity, under this jurisdiction's rules.} In
human-operated finance that mapping is checked --- an accountant, a controller,
an auditor stands somewhere in the loop. The premise of the M2M agentic economy
is that this loop closes without a human on the critical path of each
transaction: agents transact at a frequency and volume where per-transaction
human review is not merely expensive but structurally impossible.

That premise relocates a hard problem. If no human verifies each mapping, then
the \emph{trustworthiness} of the mapping can no longer live in the review
\emph{process}; it must live in the \emph{architecture} that produces the
mapping. This is the question this paper isolates for the multi-agent-systems
community. Settlement-layer protocols --- AP2 for payment
authorization~\cite{ap2_spec}, A2A for agent interoperation~\cite{a2a_spec} ---
answer ``did the value move, and was it
authorized?''\ They do not answer ``was the accounting semantics of the moved
value correct and reconstructible?''\ The latter is a property of the
\emph{semantic-resolution layer} that sits between an agent's intent and the
ledgers it touches, and it is where an early error is most dangerous: a
misresolved account at transaction time contaminates every downstream
consolidation, tax computation, and audit that inherits it.

We take Kontablo --- an existing, published, graph-based universal accounting
ontology (concept DOI 10.5281/zenodo.20738795) --- as the case study, not the
subject. The subject is the class of invariants any such layer needs. Kontablo
is useful here precisely because it already commits to determinism as a design
driver (its principle~\#5: move every decision that \emph{can} be made by a
rule, a constraint, or a graph lookup off of stochastic inference), so it lets
us make the invariants concrete and, critically, \emph{regenerable from
committed code}.

\summarynote{The claim of this paper: for autonomous high-frequency finance,
trust in the account-resolution step must be an architectural guarantee
(verifiable, reproducible), not a procedural one (human review). We give three
invariants that make it so, and a committed audit that exhibits them.}

% =====================================================================
\section{What is the resolution layer, and where can it lose information?}
% =====================================================================

A universal accounting layer projects many local statutory charts onto one
shared ontology. Kontablo keys every universal concept by a UUID and treats
local codes as visual labels (its principles \#1--\#2: a graph, not a tree;
UUID as truth). An agent asks the layer to resolve a local account --- say a
Brazilian SPED line or a German SKR04 Kontenklasse --- to a universal node so
that value booked in one jurisdiction can be reasoned about, consolidated, or
settled alongside value booked in another.

The projection is, by construction, not a bijection. Several distinct local
concepts may legitimately collapse into one coarser universal node; the return
trip cannot recover which local concept produced a given universal figure unless
the layer \emph{stored} that fact. A layer can therefore lose information in
three ways, which we must distinguish honestly (following the case study's
ADR-016):

\begin{itemize}
  \item \textbf{Lineage loss} --- the resolution is computed but the
  per-entry provenance is discarded at the API boundary, so no surface can
  answer ``which local rows produced this consolidated figure?''\ This is an
  \emph{accidental} loss and is eliminable.
  \item \textbf{Structure loss} --- the local chart's own hierarchy and its
  analytical dimensions (tax regime, VAT rate, geography, contra polarity) are
  flattened away by the mapping. Also accidental; eliminable by preserving the
  source structure alongside the projection.
  \item \textbf{Semantic coarsening} --- $N$ local concepts genuinely map to
  one deliberately coarser universal node. This is \emph{inherent}: the
  universal node is coarser on purpose. It cannot be made bijective, but it can
  be made \emph{explicit, typed, and reconstructible}.
\end{itemize}

The distinction the community should take away is: \textbf{the projection is
lossy; the system need not be.} A layer is lossless if it stores the
\emph{morphism} --- the annotated mapping plus the intact source --- rather than
only the projected image. The universal view then becomes a \emph{query over
preserved data}, never a destructive transform. The three invariants below are
the operational form of that principle for an autonomous agent.

% =====================================================================
\section{Which three invariants make the resolution step verifiable?}
% =====================================================================

\subsection{I1 --- Ontology-as-constraint: the agent's tools are deterministic,
and the graph is the guardrail}

The first invariant constrains \emph{what an agent can even ask the layer to
do}. Kontablo exposes an agent-native face over the Model Context Protocol
(MCP): a set of tools an autonomous agent calls directly. The design decision
that matters for trust is what is \emph{not} a tool.
Figure~\ref{fig:agent-flow} shows the whole decision path an agent can reach,
including the branch it cannot.

% =====================================================================
% Figure 1 — Agent -> resolution-layer decision flow (invariants I1 + I2)
%
% Structural diagram. Every label is taken verbatim from committed code;
% nothing here is illustrative-only. Provenance of each label:
%
%   6 MCP tools, none calling a model .... api/mcp/server.py
%                                          (6x "@server.tool(" decorators)
%   Tier-1 exact index lookup ............ core/harness/resolution.py:69
%                                          by_code[jurisdiction][code]
%   tier="tier1_exact", conf=1.0,
%   rule_id="tier1:<jur>:<code>" ......... core/harness/resolution.py:70-75
%   Tier-2 deterministic keyword rules ... core/harness/resolution.py:25-55
%                                          (TIER2_RULES, 10 languages)
%   tier="tier2_keyword", conf=0.85,
%   rule_id="tier2:<node>:<keyword>" ..... core/harness/resolution.py:81
%   escalation (None,"escalated",0.0,None) core/harness/resolution.py:83
%   resolved=false + escalate note ....... api/mcp/server.py resolve_account_impl
%   MappingQuote field names ............. core/harness/provenance.py:36-44
%   Tier 3 exists but is off the tool
%   surface .............................. core/harness/resolution.py:9-10
%                                          api/src/services/mapping.py
%   loss-ledger record types ............. scripts/roundtrip_audit.py ->
%                                          research/experiments/roundtrip_audit/
%                                          results.json ("loss_ledger")
%
% Regenerate/verify the labels:
%   grep -c "@server.tool(" api/mcp/server.py            # -> 6
%   sed -n '58,90p' core/harness/resolution.py           # tiers + rule_id
%   sed -n '28,49p' core/harness/provenance.py           # MappingQuote fields
%
% Requires: tikz with arrows.meta, positioning, fit, backgrounds.
% NB: `step` and `stop` are reserved pgf/TikZ keys — the style names below are
% deliberately non-colliding (tierbox / refuse / outside).
% =====================================================================
\begin{figure}[t]
\centering
\fitwidth{%
\begin{tikzpicture}[
    >=Latex,
    font=\small,
    every node/.style={inner sep=4pt},
    agent/.style={draw=clappsInk!80, rounded corners=4pt, fill=clappsInk!7,
                  text width=4.6cm, align=center, minimum height=0.95cm,
                  font=\small\bfseries},
    hdr/.style={draw=none, fill=none, text width=12.6cm, align=left,
                font=\scriptsize},
    tierbox/.style={draw=clappsAccent!80!black, rounded corners=4pt,
                    fill=clappsAccent!10, text width=4.3cm, align=center,
                    minimum height=1.0cm, font=\small},
    refuse/.style={draw=clappsAccent!80!black, rounded corners=4pt,
                   fill=clappsAccent!4, text width=4.3cm, align=center,
                   minimum height=1.0cm, font=\small},
    quote/.style={draw=clappsAccent!65!black, rounded corners=4pt,
                  fill=white, text width=5.4cm, align=left,
                  minimum height=0.95cm, font=\scriptsize},
    esc/.style={draw=red!55!black, rounded corners=4pt, fill=red!7,
                text width=5.4cm, align=left, minimum height=0.95cm,
                font=\scriptsize},
    ledger/.style={draw=red!55!black, rounded corners=4pt, fill=red!5,
                   text width=5.4cm, align=center, minimum height=0.9cm,
                   font=\scriptsize},
    outside/.style={draw=black!45, dashed, rounded corners=4pt, fill=black!5,
                    text width=5.0cm, align=center, font=\scriptsize\itshape,
                    text=black!55},
    lbl/.style={font=\scriptsize, inner sep=2pt}
]

% ---------------------------------------------------------------- agent
\node[agent] (agent) at (-4.8,3.55)
  {Autonomous agent\\[1pt]
   {\scriptsize\mdseries settles over AP2 / A2A;\\ no human in this loop}};

% ------------------------------ header strip (inside the tool surface)
\node[hdr] (hdr) at (3.0,1.95)
  {{\bfseries\color{clappsAccent!75!black}Agent-facing tool surface (MCP):
    six deterministic tools; none calls a language model.}\\[1pt]
   {\itshape Ontology-as-constraint} (\textbf{I1}) --- the identifier returned
   is always a key of the committed graph
   (\texttt{core/schemas/level3\_accounts.yaml}); there is no generative step
   that could emit one that does not exist.};

% ------------------------------------------------- deterministic spine
\node[tierbox] (t1) at (0.6,0.35)
  {\textbf{Tier 1} --- exact statutory\\ index lookup};
\node[tierbox] (t2) at (0.6,-1.85)
  {\textbf{Tier 2} --- deterministic\\ keyword rule};
\node[refuse] (esc) at (0.6,-4.05)
  {\texttt{resolved = false}\\[1pt]
   {\scriptsize refuse to guess}};

% ------------------------------------------------------------- outcomes
\node[quote] (q1) at (7.0,0.35)
  {\textbf{\texttt{MappingQuote}}\\
   \texttt{tier1\_exact} $\cdot$ conf \texttt{1.0}\\
   \texttt{rule\_id: tier1:<jur>:<code>}};
\node[quote] (q2) at (7.0,-1.85)
  {\textbf{\texttt{MappingQuote}}\\
   \texttt{tier2\_keyword} $\cdot$ conf \texttt{0.85}\\
   \texttt{rule\_id: tier2:<node>:<keyword>}};
\node[esc] (q3) at (7.0,-4.05)
  {\textbf{\texttt{MappingQuote}} --- an explicit ``I do not know''\\
   \texttt{kontablo\_id: None} $\cdot$ \texttt{tier: escalated}\\
   \texttt{confidence: 0.0} $\cdot$ \texttt{rule\_id: None}};

% ------------------------------------------------ tool-surface boundary
\begin{scope}[on background layer]
  \node[draw=clappsAccent!85!black, line width=1.1pt, rounded corners=6pt,
        fill=clappsAccent!3, fit=(hdr)(t1)(t2)(esc)(q1)(q2)(q3),
        inner sep=10pt] (surface) {};
\end{scope}

% --------------------------------------------------------------- edges
% The agent's call enters the surface from outside: elbow down, then in.
\draw[->, thick] (agent.south) -- (-4.8,0.35) -- (t1.west);
\node[lbl, anchor=east, align=right] at (-4.95,1.75)
  {\texttt{resolve\_account(}\\ \texttt{code, name,}\\ \texttt{jurisdiction)}};

\draw[->, thick] (t1) -- node[lbl, right] {no match} (t2);
\draw[->, thick] (t2) -- node[lbl, right] {no match} (esc);
\draw[->, thick] (t1) -- node[lbl, above] {hit} (q1);
\draw[->, thick] (t2) -- node[lbl, above] {hit} (q2);
\draw[->, thick] (esc) -- (q3);

% ------------------------------------------- Tier 3: outside the surface
\node[outside] (t3) at (0.6,-7.55)
  {Tier 3 --- semantic / LLM fallback\\[2pt]
   \upshape exists in the codebase, \textbf{not} registered as a tool};
\draw[->, dashed, black!45] (esc) -- (t3);
\node[font=\normalsize\bfseries, text=red!65!black] at (0.6,-6.05) {$\times$};
\node[lbl, anchor=east, text=red!65!black, align=right] at (0.3,-6.05)
  {never reachable from the\\ agent-facing tool surface};

% ------------------------------------------------ loss ledger + human
\node[ledger] (ll) at (7.0,-6.25)
  {\textbf{Loss ledger (I2)} --- typed and counted:\\
   collisions $\cdot$ placeholders $\cdot$ escalations $\cdot$ CRA flags\\
   \texttt{silent\_losses = 0}};
\node[ledger, fill=white, draw=clappsInk!55] (human) at (7.0,-8.35)
  {Human review: the human remains the legal
   principal, consulted where certainty runs out
   --- not on every transaction};
\draw[->, thick, red!55!black] (q3) -- (ll);
\draw[->, thick, clappsInk!60] (ll) -- (human);

\end{tikzpicture}%
}
\caption{How an autonomous agent reaches the resolution layer, and what it can
get back. The agent's entire callable surface is deterministic: Tier~1 is an
exact lookup into the committed statutory index, Tier~2 a fixed multilingual
keyword rule, and every outcome --- including the refusal --- is stamped with a
\texttt{MappingQuote} naming the exact rule that fired
(\protect\pathtt{core/harness/provenance.py}). \textbf{I1} is the property that
the identifier returned is always a key of the committed graph, so a booking to
a non-existent account is unreachable rather than merely unlikely. \textbf{I2}
is the property that the third branch is a \emph{typed, counted} record rather
than a guess or a silent drop. The semantic (Tier~3) fallback exists in the
codebase but is deliberately \emph{not} registered as an MCP tool: exposing it
would let an agent consume a stochastic answer through an interface that looks
deterministic. Confidence values and \texttt{rule\_id} formats are verbatim from
\protect\pathtt{core/harness/resolution.py}.}
\label{fig:agent-flow}
\end{figure}


The reference MCP server (\pathtt{api/mcp/server.py}) registers six tools:
\texttt{resolve\_account}, \texttt{get\_account}, \texttt{validate\_balance\_sheet},
\texttt{consolidate\_trial\_balances}, \texttt{get\_node\_fiber}, and
\texttt{list\_jurisdictions}. Every one is a graph lookup, a deterministic
keyword rule, or arithmetic --- \emph{no tool calls a language model}. The
resolver has a three-tier structure (an exact statutory index, a deterministic
keyword rule, and a semantic/LLM fallback), and the third tier is
\textbf{deliberately not exposed as a tool}. When neither deterministic tier
matches, \texttt{resolve\_account} returns \texttt{resolved=false} and escalates
to human review; it never guesses. Exposing the stochastic fallback as if it
were just another tool would hide non-determinism from an agent that is deciding
autonomously --- precisely the failure this invariant exists to prevent.

The guardrail is the ontology itself: an agent cannot propose an account UUID
that does not exist in the graph, because resolution is a lookup into the graph,
not a generation of free text. This is ontology-as-constraint, and it is the
architectural (not procedural) source of the trust: the space of possible
outputs is bounded by committed data, so a whole class of hallucinated or
malformed bookings is unreachable rather than merely unlikely. The principle is
the one that constrained decoding and identifier-constrained retrieval
establish~\cite{koo2024automata,decao2021genre}, carried to the point where the
identifier-producing step has no generative surface at all; we return to that
distinction, and to its known critiques, in Section~\ref{sec:related}.

\subsection{I2 --- The loss ledger: nothing untranslatable is ever silent}

The second invariant governs \emph{what happens at the boundary of the layer's
competence}. The standing rule is: \emph{the pipeline may fail to translate, but
it may never lose.} Every non-translation is a typed, countable record ---
an ontology code collision, a non-code placeholder, an escalated entry, a
compliance (co-responsibility) flag --- never a silent drop. We call the running
tally the \emph{loss ledger}, and its load-bearing quantity is
$\texttt{silent\_losses}$: information discarded \emph{without} a typed record.
The invariant is $\texttt{silent\_losses} = 0$.

Concretely, every resolution carries a \texttt{MappingQuote}
(\pathtt{core/harness/provenance.py}) --- the analogue of a per-transaction FX
quote --- recording the local code and name, the jurisdiction, the tier that
answered, the \emph{exact deterministic rule identifier} that fired
(\texttt{tier1:<jur>:<code>} or \texttt{tier2:<node>:<keyword>}), a confidence,
and a timestamp. Escalations carry a \texttt{MappingQuote} with a null node and
tier \texttt{escalated} --- an explicit ``I do not know,'' not an absence.
Because the quote is attached to every resolved entry and surfaced on every
face (REST, gRPC, MCP), a consolidated statement is auditable down to the
individual booking that produced each figure.

This is a stronger commitment than auditability as the surrounding literature
usually defines it. Frameworks for auditable agents are concerned with whether a
past action can be recovered and inspected~\cite{nian2026auditableagents,wang2026agenttraces};
the loss ledger additionally requires that the \emph{absence} of a translation
be a first-class, counted object.

For an autonomous agent this is the difference between an error it can detect and
an error it cannot. A silent loss is invisible until a much later reconciliation
fails, by which point the contaminated fact has propagated through every agent
that consumed it. A ledgered loss is a signal the agent can act on \emph{at the
boundary} --- pause, escalate, seek a human --- keeping the human in the loop
where certainty runs out, without putting a human on the critical path where it
does not.

\subsection{I3 --- The pre-transaction fiber query: audit before you act}

The third invariant concerns \emph{when} auditability is available. Post-hoc
audit is standard: after the fact, reconstruct what happened. It is insufficient
for an agent about to commit an irreversible transaction, because the damage is
already done by the time a forensic pass runs. The invariant is that the
preimage of any node is queryable \emph{before} the agent acts on it.

Kontablo exposes the \emph{fiber} of a universal node --- the set of local
statutory codes that collapse into it, per jurisdiction, enriched with the
preserved local structure (local parent, analytical facets, declared
aggregation group). It is available as the MCP tool \texttt{get\_node\_fiber}
and as REST \texttt{GET /accounts/\{id\}/fiber}, both over the same
deterministic function \texttt{node\_fiber} (\pathtt{core/harness/ontology.py}).
Figure~\ref{fig:fiber-query} shows an actual fiber, generated from the committed
ontology rather than drawn: the universal cash node stands for $57$ distinct
local statutory codes across $56$ jurisdictions, and the query returns, per
jurisdiction, exactly which ones --- with the local parent and analytical facets
that the projection would otherwise flatten away.
Dually, \texttt{rollup(accounts, lens)} partitions the same UUIDs under
different lenses (e.g.\ IFRS presentation vs.\ cash-flow), realizing the ``one
UUID, many parallel hierarchies'' that the graph principle promises. Nodes that
do not carry a lens are grouped explicitly under a null key --- never silently
dropped, consistent with I2.

The agentic consequence: before an agent books value against a universal node,
it can ask ``exactly which local concepts does this node stand for in this
jurisdiction, and what structure was flattened to reach it?''\ and get a
deterministic answer. Auditability becomes a \emph{precondition} of the
transaction, not a reconstruction after it. In a high-frequency M2M setting this
is what lets an agent make a bounded, checkable decision at machine speed:
the preimage is a fact it can inspect and reason over, not a promise it must
trust.

Two precisions keep this claim the size it should be. \emph{First, the
mathematics is elementary and we do not dress it up.} Write $\pi$ for the
projection from the disjoint union of local statutory codes
$\coprod_{j} C_{j}$ onto the universal node set $V$; the fiber of a node
$v \in V$ is the ordinary preimage $\pi^{-1}(v)$, and its restriction to one
jurisdiction is $\pi^{-1}(v) \cap C_{j}$. We use ``fiber'' in exactly that
elementary sense --- no bundle, no local triviality, no sheaf structure is
asserted here. That many local codes collapse onto one universal node is not a
discovery; it is what a coarsening \emph{is}.

What is claimed is availability and content, not existence. \emph{Availability:}
the preimage is materialized behind the same deterministic surface the agent
uses to post, so it is answerable before the transaction rather than
reconstructible after it. \emph{Content:} the query returns not the bare set
$\pi^{-1}(v)$ but the \emph{annotated} preimage --- each member carrying the
local parent, the analytical facets and the declared aggregation group that the
projection discards, plus a \texttt{source} tag distinguishing a Tier-1 index
member from a localization member. That is the operational form of storing the
morphism rather than the image (Section~2): it makes the return trip a query
over preserved data instead of an inference about lost data. A standard that
publishes a periodic export can also, in principle, be inverted; what it cannot
do is answer the question at the moment the agent needs it, over the interface
the agent is already holding.

\emph{Second, we stop short of the formal account on purpose.} The natural
question behind I3 --- when can the local ledgers of many jurisdictions be glued
into one globally consistent whole, and what obstructs it? --- is a
cohomological one, and the hub sketches the intended referent: a cellular sheaf
over the mapping graph whose $H^{0}$ is the space of globally consistent
reconciliations and whose non-vanishing $H^{1}$ localizes the obstruction to
gluing~\cite{kontablo_hub}. The hub is explicit that this is an \emph{intended}
formal referent, validated by a committed toy experiment rather than proved in
general, and we inherit that caveat rather than laundering it into a stronger
claim by restating it. I3 is a systems result: the fiber query supplies exactly
the data such an account requires --- the local sections, intact, per
jurisdiction --- without asserting the theory that would consume them. The
formal development belongs to a forthcoming companion spoke on the mathematical
foundations, which this paper deliberately does not pre-empt. Everything claimed
here stands or falls on the committed audit of Section~4, not on the cohomology.

%% =====================================================================
%% Figure 2 --- the pre-transaction fiber query (invariant I3)
%%
%% GENERATED FILE --- do not edit by hand.
%% Regenerate with:  python docs/papers/spokes/agentic-provenance/figures/gen_fig_fiber_query.py
%%
%% Every code, name, local_parent and source tag below is emitted by
%% core.harness.ontology.node_fiber() over the committed ontology and
%% localization files --- the same function behind the MCP tool
%% get_node_fiber and REST GET /accounts/{id}/fiber.
%% =====================================================================
\begin{figure}[t]
\centering
\fitwidth{%
\begin{tikzpicture}[
    >=Latex,
    font=\small,
    ask/.style={draw=clappsInk!80, rounded corners=4pt, fill=clappsInk!7,
                text width=8.6cm, align=center, minimum height=0.9cm,
                font=\small},
    uni/.style={draw=clappsAccent!85!black, line width=1.0pt,
                rounded corners=4pt, fill=clappsAccent!10, text width=8.6cm,
                align=center, minimum height=1.1cm, font=\small},
    fiber/.style={draw=clappsAccent!60!black, rounded corners=4pt, fill=white,
                  text width=4.55cm, align=left, font=\scriptsize,
                  anchor=north},
    lbl/.style={font=\scriptsize, inner sep=2pt}
]

\node[ask] (ask) at (0,2.05)
  {\textbf{Agent, \emph{before} committing the transaction:}\\[2pt]
   \texttt{get\_node\_fiber("asset.current.cash", jurisdiction)}};

\node[uni] (node) at (0,-0.45)
  {\textbf{universal node} \texttt{asset.current.cash}\\[1pt]
   {\scriptsize \texttt{00000000-0000-4000-8000-000000000101} --- ``Cash and Cash Equivalents''}\\[2pt]
   {\scriptsize preimage: 57 local codes across
    56 jurisdictions (DE, BR, MX shown)}};

\draw[->, thick, clappsInk!70] (ask) -- (node);

\node[fiber] (de) at (-5.35,-2.60)
  {\textbf{DE --- SKR (Germany)}\\[3pt] \texttt{1000}\\[-1pt] {\scriptsize {\itshape tier1\_index}}\\[3pt] \texttt{1600} --- Kasse\\[-1pt] {\scriptsize parent \texttt{1} $\cdot$ {\itshape localization}}};
\draw[->, thick, clappsAccent!70!black] (node.south) -- (de.north);
\node[fiber] (br) at (0.00,-2.60)
  {\textbf{BR --- SPED (Brazil)}\\[3pt] \texttt{1.01.01.01} --- Caixa Geral\\[-1pt] {\scriptsize parent \texttt{1.01.01} $\cdot$ {\itshape localization}}\\[3pt] \texttt{1.1.1.1.01-9}\\[-1pt] {\scriptsize {\itshape tier1\_index}}};
\draw[->, thick, clappsAccent!70!black] (node.south) -- (br.north);
\node[fiber] (mx) at (5.35,-2.60)
  {\textbf{MX --- SAT (Mexico)}\\[3pt] \texttt{101} --- Caja\\[-1pt] {\scriptsize parent \texttt{100} $\cdot$ {\itshape tier1\_index}}};
\draw[->, thick, clappsAccent!70!black] (node.south) -- (mx.north);

\end{tikzpicture}%
}
\caption{The pre-transaction fiber query (\textbf{I3}). Given a universal node,
the layer returns its \emph{preimage}: which local statutory codes collapse into
it in each jurisdiction, with the local structure the projection would otherwise
flatten away (the local parent, the analytical facets, the declared aggregation
group), and a \texttt{source} tag saying whether the member came from the
deterministic Tier-1 index or from the jurisdiction's localization file. Because
the query is available \emph{before} the agent posts, auditability is a
precondition of the transaction rather than a forensic reconstruction after it.
This figure is generated, not drawn: every code and label is emitted by
\texttt{node\_fiber} (\protect\pathtt{core/harness/ontology.py}) over the committed
ontology, via \protect\pathtt{docs/papers/spokes/agentic-provenance/figures/gen_fig_fiber_query.py}.}
\label{fig:fiber-query}
\end{figure}


\subsection{What can a hostile or buggy agent do to these invariants?}
\label{sec:threat}

The invariants are only interesting if they survive an agent that is not
cooperative. We therefore state the threat model explicitly. The trusted
computing base is the committed ontology, the resolver code, and the
localization data; the \emph{agent is untrusted}. We assume it can call any
exposed tool, arbitrarily often, with arbitrary arguments --- whether because it
is adversarial, miscalibrated, or simply buggy. The agent supplies three fields
to \texttt{resolve\_account}: a jurisdiction, a local code, and a local name.

Three things it cannot do. It cannot cause the layer to emit an identifier that
does not exist (I1): resolution is a dictionary lookup whose value is by
construction a key of the loaded account graph, so there is no path from
agent-controlled input to a fabricated UUID. It cannot cause a silent loss (I2):
the worst outcome of garbage input is a \emph{counted} escalation, which is a
record, not a disappearance --- an agent can make the loss ledger larger, never
emptier. And it cannot hide a preimage (I3): \texttt{node\_fiber} is a read over
committed data, and no tool on the surface mutates the ontology.

What it \emph{can} do is where the honesty matters.

\begin{itemize}
  \item \textbf{Escalation flooding.} An agent can submit unresolvable entries
  at machine speed and inflate the escalation queue until human review capacity
  is exhausted. I2 guarantees each such entry is typed and counted; it does not
  guarantee the queue stays tractable. That is an operational control ---
  per-principal rate limits and quotas --- and the reference implementation does
  not have one. We claim the accounting property, not availability.

  \item \textbf{A wrong but existing node.} This is the decisive residual. I1
  bounds the output space to real nodes; it does not make the chosen node the
  \emph{correct} one. Tier-2 matches deterministic keyword rules against the
  agent-supplied \emph{name}, so an agent that controls that string can steer
  resolution toward a plausible-but-wrong node --- a name containing
  ``cash'' that denotes something else will resolve to the cash node. Tier~1 is
  not steerable by naming, since it keys on (jurisdiction, code) against the
  committed index, but an agent that asserts the wrong local code will be
  faithfully resolved to whatever that code means. The invariants govern the
  \emph{translation}; they do not certify the truthfulness of the assertion
  being translated.

  What I1 and I2 do buy against this case is attribution. The wrong answer is a
  real node accompanied by a \texttt{MappingQuote} naming the exact rule that
  fired --- \texttt{tier2:<node>:<keyword>} identifies both the target and the
  substring that produced it --- so the error is reviewable by querying
  resolutions of one tier, and reproducible without re-running the resolver. The
  contrast is with a generated identifier, where there is nothing to query. The
  failure class is converted from unattributable to attributable, not
  eliminated; this is the same ``bounded, not eliminated'' posture I1 takes
  throughout.

  \item \textbf{A deterministic boundary library that is itself
  name-dependent.} A partial defence does exist:
  \texttt{cra\_validate} (\pathtt{core/harness/boundary.py}) applies
  accounting invariants that can catch a mis-steer --- a declared debit/credit
  nature that disagrees with the target node, and a liquid-asset name landing on
  a non-current node, both checked on every resolution; statement-class, VAT
  direction and equity-versus-liability confusions, checked on proposed
  (forced) mappings. But these checks test \emph{consistency between the
  supplied name and the target}. An agent that supplies a name and a code that
  are mutually consistent and jointly wrong passes all of them. We report this
  as a bound on the boundary library, not as a gap to be papered over.

  \item \textbf{Confidence is not correctness.} A \texttt{confidence} of
  \texttt{1.0} on a Tier-1 hit means an exact index entry existed, not that the
  booking is economically right. A downstream agent that reads the field as a
  probability of correctness is misreading it, and the field name invites that
  misreading.

  \item \textbf{No authentication here.} The layer does not establish who the
  calling agent is or whether it is authorized to post. Identity, mandate and
  authorization belong to the settlement layer (AP2's signed mandates,
  ERC-8004's on-chain identity), and MCP's own specification states it ``cannot
  enforce these security principles at the protocol
  level''~\cite{mcp_spec}. These invariants assume authentication; they do not
  supply it.
\end{itemize}

The summary a reviewer should take: the invariants relocate the locus of error
to the semantic coverage boundary --- was the right concept named? --- and make
everything on the near side of that boundary verifiable, counted, and
attributable. They do not claim the boundary itself has been abolished.

% =====================================================================
\section{Result: a round-trip audit with zero silent losses}
% =====================================================================

The invariants would be marketing if they were not regenerable. Kontablo's
governing scientific rule is that \emph{every quantitative claim in a public
artifact must be reproducible from a committed, deterministic command.} We honor
it here: every number below is emitted by \pathtt{scripts/roundtrip_audit.py}
into \pathtt{research/experiments/roundtrip_audit/results.json}, and the same
script runs in CI as a hard gate.

The audit runs on the \emph{same} synthetic dataset as the project's published
consolidation validation (imported, not duplicated), so no rows are constructed
for this paper. The two scripts do report different entry totals, for a
deliberate reason worth stating plainly: the coverage run counts $366$ entries
because it excludes the single per-entity balancing (equity-plug) line, so that
trivially-resolving plugs cannot pad the deterministic-coverage percentage; the
round-trip audit keeps those same lines as ordinary ledger rows, because its
question is conservation rather than coverage, giving $441 = 366 + 75$ (one plug
per entity). Both figures are correct for what they measure. On that dataset the
audit verifies four properties:

\begin{enumerate}
  \item \textbf{Conservation.} Every input trial-balance row appears in the
  engine's resolved output exactly once --- $441$ entries in, $441$ resolved
  out; zero missing rows, zero phantom rows.
  \item \textbf{Reconstruction.} Every local trial balance (code, name, side,
  local-currency amount) is rebuilt byte-for-byte from the lineage alone,
  across all $75$ entities.
  \item \textbf{Fiber consistency.} Every consolidated line equals the sum of
  its fiber --- the resolved entries that aggregated into it --- with zero
  mismatches.
  \item \textbf{Loss ledger.} Everything the pipeline could not translate is a
  typed record, and the count of \emph{silent} losses is zero.
\end{enumerate}

\begin{table}[H]
\centering
\small
\begin{tabular}{@{}lr@{}}
\toprule
\textbf{Round-trip audit (ADR-016)} & \textbf{Value} \\
\midrule
Entities (synthetic) & 75 \\
Jurisdictions & 68 \\
Entries in $\rightarrow$ resolved out & $441 \rightarrow 441$ \\
Missing rows / phantom rows & $0$ / $0$ \\
Reconstruction exact & true \\
Fiber mismatches & $0$ \\
Provenance complete (every entry) & true \\
\midrule
\multicolumn{2}{@{}l}{\emph{Loss ledger} (typed, countable --- not silent):}\\
\quad Ontology code collisions & 0 \\
\quad Non-code placeholders & 14 \\
\quad Escalated entries & 4 \\
\quad Co-responsibility (CRA) flags & 2 \\
\midrule
\textbf{Silent losses} & \textbf{0} \\
\bottomrule
\end{tabular}
\caption{Output of \texttt{python scripts/roundtrip\_audit.py} (exit code 0).
The four escalated entries are exactly those the resolver declines to map
deterministically; the count matches the published $97.3\%$ run
(\texttt{escalated\_entries: 4} in both artifacts), and since both scripts drive
the identical resolver over the identical generated dataset, they are expected
to name the same four rows. Escalations, placeholders, collisions and compliance
flags are the loss ledger; they are explicit and counted, which is why the
silent-loss count is zero.}
\label{tab:roundtrip}
\end{table}

The reading for the agent community: the loss ledger makes the difference
between ``the system handled everything'' and ``the system handled everything
\emph{it was certain about}, and told you exactly what it was not certain
about.'' The second is the honest and the safe posture for an autonomous
transaction pipeline. The $14$ placeholders and $4$ escalations are not the
system failing; they are the system refusing to guess, in a form an agent can
detect and route.

\subsection{What does moving the decision off inference buy?}
\label{sec:inference-cost}

One consequence of I1 deserves stating, because it is easy to overclaim and we
want the modest version on the record. In the reference implementation the
resolution step makes \emph{no} inference calls at all: Tier-3 is not on the
tool surface, so the count is zero by construction rather than by measurement.
The interesting quantity is therefore counterfactual --- what a layer of this
shape would spend if it \emph{did} fall back to a model on the residual.

On the round-trip dataset, $4$ of $441$ entries escalate ($0.9\%$); the other
$437$ ($99.1\%$) are settled by an index lookup or a fixed keyword rule. In the
coverage run the same resolver reports $310$ Tier-1 and $46$ Tier-2 resolutions
against $366$ counted entries, $97.3\%$ deterministic. So a layer that routed
its residual to a language model would issue inference calls proportional to the
\emph{escalation rate}, not to the transaction count. That is the structural
point: determinism at the front makes the inference budget a function of the
semantic coverage boundary rather than of transaction volume --- which is the
difference that matters when the premise of the setting is machine-frequency
transaction rates.

The same asymmetry applies to error, and this is the more important of the two.
A stochastic decision taken early contaminates every step that inherits it: a
misresolved account propagates into consolidation, tax computation, and audit.
A deterministic decision taken early is inherited cleanly. Bounding inference to
the residual bounds the contaminable surface to the residual as well.

We deliberately do not attach token, latency or energy figures to any of this.
We have run no inference-cost benchmark, and a call that is never made is not a
measured saving --- it is an absent cost. The claim is structural and the
escalation rate that grounds it is a property of \emph{this} synthetic dataset
against the current ontology coverage, not a constant: a real multi-jurisdiction
ledger with more idiosyncratic local naming should be expected to escalate more
often, which moves the same argument without changing its shape.

% =====================================================================
\section{Related work: where these invariants sit}
\label{sec:related}
% =====================================================================

Four bodies of work bound this contribution: the protocols that move value
between agents, the techniques that bound what a model can emit, the multi-agent
literature on holding agents accountable, and the standards that encode
financial semantics for machines. We take each in turn, and state at the end
what is left over.

\subsection{Agent interoperation, payment, and settlement protocols}

The agent-commerce stack matured rapidly over 2024--2026, and it is layered:
MCP standardizes an agent's access to tools and resources~\cite{mcp_spec}; A2A
standardizes agent-to-agent task delegation~\cite{a2a_spec}; AP2 adds payment
authorization via signed mandates~\cite{ap2_spec}; UCP~\cite{ucp_spec} and
ACP~\cite{acp_spec} orchestrate commerce on top of them --- UCP over four
declared transports including MCP and A2A, ACP over its own HTTP/OpenAPI
checkout surface with an MCP transport binding under open proposal;
x402~\cite{x402_spec} and ERC-8004~\cite{erc8004} add settlement and on-chain
identity and reputation.

What matters here is what these specifications say about the \emph{accounting
semantics} of the value they move: nothing, and they say so themselves.
ERC-8004 states that ``payments are orthogonal to this protocol and not covered
here''~\cite{erc8004}; x402 carries an explicit out-of-scope
list~\cite{x402_spec}; AP2 places the details of the commerce protocol outside
its own scope~\cite{ap2_spec}; and MCP's specification is explicit that MCP
``cannot enforce these security principles at the protocol
level''~\cite{mcp_spec}. Across the eleven protocols and vendor programs
surveyed for this paper, none defines, validates, or mentions which economic
concept a moved value should be posted to. They converge instead on four
concerns: identity, authorization, value transfer, and interoperation. The
eleven are enumerated, with their layer, governance and the date each fact was
verified against its primary specification, in the case study's committed
protocol register (\pathtt{docs/strategy/protocol-compatibility-register.md});
that file is maintained as the landscape moves, so the survey's scope is
checkable rather than asserted, and a reader can see when a claim here was last
confirmed.

This is not a deficiency in those protocols --- it is a layer boundary, and the
surrounding literature reads it the same way. A recent systematization of
blockchain agent-to-agent payments names ``limited accountability'' among the
open problems of agent payment rails~\cite{zhang2026sok}, and a comparative
study of inter-agent trust models across A2A, AP2, and ERC-8004 finds six trust
primitives in use --- brief, claim, proof, stake, reputation, and constraint ---
none of which asks whether the posted value is semantically
correct~\cite{hu2025interagenttrust}. The settlement layer knows it has an
accountability gap; its proposed fixes are cryptographic and reputational, which
is orthogonal to, not competing with, semantic correctness.

\subsection{Bounding the output space: constrained generation and ontologies}

I1 belongs to an established line of work on making invalid output
\emph{unreachable} rather than merely unlikely. Grammar- and
automaton-constrained decoding masks token logits so that sequences outside a
formal language carry zero probability
mass~\cite{willard2023outlines,geng2023gcd}; Koo et al.\ give an explicit
correctness proof that decoding cannot leave the
constraint~\cite{koo2024automata}. This is categorically different from
statistical mitigation, which shifts a distribution without bounding it.

The closest structural analogues to resolving a local code onto a UUID that must
already exist are not grammar papers but identifier-constrained ones. GENRE
constrains autoregressive generation with a prefix trie built from a knowledge
base, so the model can only ever complete a real entity
identifier~\cite{decao2021genre}; Graph-constrained Reasoning does the analogous
thing over knowledge-graph paths and reports ``accurate reasoning with zero
reasoning hallucination''~\cite{luo2024gcr}. PICARD is the classic precedent in
this shape, rejecting invalid table and column names during
decoding~\cite{scholak2021picard}.

One precision is owed to this literature. All of the above constrain a language
model's \emph{own} decoding. The Tier-1 and Tier-2 resolvers described here do
no decoding at all --- they are deterministic lookups over committed data, and
the one tier that would involve a model is deliberately absent from the tool
surface. I1 is therefore not an instance of constrained decoding in the
technical sense those papers use the term; it is the same principle taken to its
limit, with \emph{zero} generative surface for the identifier rather than a
masked one. That distinction also disposes of the strongest critique of this
family of techniques: constraining a model's reasoning tokens to a rigid grammar
measurably degrades accuracy on reasoning-heavy
tasks~\cite{tam2024letmespeak,banerjee2025crane}. That critique lands on systems
that constrain reasoning; it has no purchase on a lookup that does none, which
is precisely why the stochastic tier is kept off the tool surface rather than
wrapped in a grammar.

The empirical stakes for preferring a constructive guarantee are documented in
an adjacent high-stakes domain. Commercial, retrieval-grounded legal research
tools marketed as reliable still hallucinate 17--33\% of the
time~\cite{magesh2024hallucinationfree}, and general-purpose models hallucinate
on 58--88\% of verifiable legal queries~\cite{dahl2024legalfictions}. Grounding
lowers a rate; it does not close a class.

\subsection{Accountability, provenance, and auditability in multi-agent systems}

The multi-agent community has asked how to hold autonomous systems responsible
at least since the Blue Sky agenda of Yazdanpanah et
al.~\cite{yazdanpanah2021responsibility}. The recent literature answers
overwhelmingly in terms of \emph{reconstruction}. ``Auditable Agents'' defines
five dimensions of auditability --- action recoverability, lifecycle coverage,
policy checkability, responsibility attribution, and evidence integrity --- all
concerning whether a past action can be recovered and
inspected~\cite{nian2026auditableagents}. A survey of evidence tracing frames
the field as execution provenance: a typed graph of an agent's execution and its
projection onto support relations~\cite{wang2026agenttraces}. Atomix supplies
transactional semantics for tool use so that effects settle together and partial
failures do not leak~\cite{mohammadi2026atomix} --- the closest infrastructure
analogue to I2, but over generic tool effects rather than typed accounting
entries, and with no notion of an entry the system declines to translate.

The closest architectural cousin is AuditFlow, which builds an executable
symbolic environment from a US-GAAP taxonomy and an XBRL filing graph, exposes
it to agents as typed tools, and reports accuracy falling from 82.09\% to
17.91\% when the deterministic check layer is ablated~\cite{wang2026auditflow}.
That ablation is strong independent evidence for the general thesis that
symbolic grounding carries the weight in financial agent tasks. The difference
is position in the transaction lifecycle: AuditFlow \emph{verifies already-filed
reports} --- a detection task over static documents --- whereas I1 constrains
what an agent may post \emph{at transaction time}, and I2 ledgers what that
constraint cannot resolve. Detection after filing and prevention at posting are
complementary, and a reader who knows AuditFlow should read this paper as
addressing the other end of the same pipeline.

Stating the gap carefully: across the venues and papers surveyed here we did not
find work treating the ontology itself as the mechanism that makes a class of
posting error unreachable at authoring time, nor a named, typed, continuously
gated \texttt{silent\_losses} quantity. A broad recent agenda for multi-agent
security likewise taxonomizes threats without naming accounting-semantic
correctness as a threat class~\cite{schroederdewitt2025masopenchallenges}. We
report this as an absence within a bounded survey --- 2024--2026,
English-language, across AAMAS, AAAI, ICLR, ICML, and arXiv --- not as a claim
that no such work exists anywhere.

\subsection{Machine-facing financial semantics standards}

Kontablo is not the first attempt to give machines financial semantics, and the
comparison must be made carefully, because a claim resting on granularity alone
would be wrong. FIBO models instruments, contracts, and legal entities in
OWL/RDF, above the posting level, with no journal-entry primitives~\cite{fibo}.
XBRL~\cite{xbrl21} and its inline profile under ESEF~\cite{esef_rts} tag
period-end reporting facts and are legally mandated across the EU and the US ---
a regulatory adoption this work does not approach and should not pretend to.
ISO 20022 standardizes the \emph{message} that moves value between institutions,
not how either side's ledger books it~\cite{iso20022}.

Crucially, three standards \emph{do} already reach journal-entry granularity:
XBRL GL~\cite{xbrl_gl}, the AICPA Audit Data Standards~\cite{aicpa_ads}, and
OECD SAF-T~\cite{oecd_saft}. Entry-level granularity is therefore not this
work's distinguishing property. What separates the resolution layer described
here is a combination none of them provides: a live, agent-callable
pre-transaction query surface, where all three are periodic batch exports
consumed after the fact; one canonical ontology spanning many sovereign
jurisdictions' charts rather than one dialect per country; per-entry lineage
sufficient to reconstruct the source; and a typed guarantee that untranslatable
material is counted rather than dropped.

The cross-jurisdiction point has an unusually clean primary-source witness.
OECD's own SAF-T guidance leaves the XML representation to individual revenue
bodies~\cite{oecd_saft}, and the predictable result is a family of mutually
incompatible national dialects --- Portugal's SAF-T PT, Poland's JPK, Romania's
D406, Norway's SAF-T Financial. The fragmentation this work targets is not
hypothetical; it is visible in the design of the standard that comes closest to
addressing it.

\subsection{Positioning, and honesty about implementation status}

We place the contribution deliberately narrowly. AP2 and A2A operate at the
settlement and interoperation layers: they authorize and move value between
agents. This paper's invariants operate one layer beneath, on the ledger's
\emph{meaning}: they make the account-resolution step verifiable and
reconstructible. The two are complementary --- a payment can be perfectly
authorized under AP2 and still be booked to the wrong universal concept if the
resolution layer guessed silently. Conversely, a lossless, ledgered resolution
layer does not settle anything; it makes the thing being settled auditable.

That layer separation has a practical consequence worth stating, because a
reviewer will reasonably ask what happens to this contribution when the protocol
landscape shifts --- and in 2026 it is shifting fast enough that no one can name
tomorrow's standards. If the invariants live beneath the protocol layer, then
protocol churn is absorbed as a \emph{transport} concern rather than a redesign:
a new face is an adapter over the same deterministic resolver, and no adapter
carries accounting logic of its own. The current stack supplies evidence that
this boundary is real and cheap to cross. Two independently developed commerce
protocols both route to the same tool transport --- UCP declares MCP among its
four transports, and ACP has an MCP transport binding under open
proposal~\cite{ucp_spec,acp_spec} --- so a single deterministic tool surface is
reachable from both without either protocol knowing anything about charts of
accounts. Breadth of compatibility here is not a feature list; it is what keeps a
semantic layer alive across a protocol generation.

We can therefore state the condition under which this positioning would fail,
which is the honest way to hold it. The claim that these invariants belong to a
distinct layer rests on an empirical observation, not a theorem: across the
protocols surveyed here, none defines, validates, or mentions which economic
concept a moved value should be posted to. Were a settlement or commerce protocol
to begin specifying accounting semantics itself --- binding a taxonomy, or
requiring that a booking be reconstructible from the payload --- the layer
boundary this paper argues for would collapse into that protocol, and the
contribution would have to be re-stated as a proposal \emph{to} it rather than a
layer beneath it. We are not aware of any protocol doing so today, and we name
the condition rather than assume it away.

Honesty about implementation status matters to this community, which has seen
much asserted and less shipped. In the case study, the MCP deterministic core is
\emph{implemented} (the six tools above) and the round-trip audit is
\emph{CI-gated}. The AP2 and A2A integrations are \emph{asserted architecture,
not running code}; the reference implementation is prototype-grade relative to
the specification. We claim the invariants and the audit that exhibits them ---
not a production settlement stack.

% =====================================================================
\section{Limitations and threats to validity}
% =====================================================================

\begin{itemize}
  \item \textbf{Synthetic data.} The $75$ entities and $441$ entries are
  synthetic trial balances, not exports from live ledgers. The audit
  demonstrates a \emph{property of the architecture} (conservation,
  reconstruction, ledgered loss), not a measurement of real-world posting
  distributions. Validating the invariants on real multi-jurisdiction ledgers
  is future work.
  \item \textbf{Semantic coarsening is still lossy.} I2 and I3 make the
  coarsening explicit and reconstructible; they do not make it bijective. Some
  formal claims of the case study's ontology (certain analytical-facet families
  and consolidation roles) are only partially closed by the current schema and
  are tracked for a later specification version.
  \item \textbf{Scope of the agent-native face.} Only MCP's deterministic core
  is implemented; the stochastic (Tier-3) resolution is intentionally absent
  from the tool surface, and A2A/AP2 are not implemented. Nothing here should be
  read as a claim of full protocol coverage.
  \item \textbf{Residual attack surface.} Section~\ref{sec:threat} states what a
  hostile or buggy agent can still achieve --- escalation flooding, and steering
  Tier-2 toward a wrong-but-existing node via the supplied account name --- and
  why neither is closed by these invariants. Authentication and rate limiting
  are assumed, not provided.
  \item \textbf{Single implementation.} The invariants are demonstrated in one
  reference system. Their generality as a design pattern for other
  resolution layers is argued, not independently replicated.
\end{itemize}

% =====================================================================
\section{Conclusion}
% =====================================================================

For autonomous agents transacting at machine frequency, trust in the
account-resolution step cannot be purchased with human review, because there is
no human in the per-transaction loop to buy it from. It has to be built into the
architecture. We named three invariants that do so --- ontology-as-constraint
(I1), the loss ledger (I2), and the pre-transaction fiber query (I3) --- and
exhibited them in a committed reference implementation whose round-trip audit
records zero silent losses over $441$ entries, $75$ entities, and $68$
jurisdictions, reconstructing every local ledger from lineage alone. The
takeaway for the multi-agent community is a concrete, regenerable shape for
machine-verifiable trust at the resolution layer: store the morphism, never the
image alone; expose only deterministic tools and let the ontology be the
guardrail; and make ``I cannot translate this with certainty'' a typed,
countable, escalating event rather than a silent guess. Anything a system cannot
prove, it should be architected to \emph{declare} --- and to declare it before
the agent acts, not after.

% =====================================================================
\section*{Reproducibility}
% =====================================================================

Every quantitative claim in this paper regenerates from the following committed
commands in the Kontablo repository (\texttt{accounting-esperanto}):

\begin{itemize}
  \item Table~\ref{tab:roundtrip} (all round-trip / loss-ledger numbers):\\
  \texttt{python scripts/roundtrip\_audit.py} $\rightarrow$
  \pathtt{research/experiments/roundtrip_audit/results.json} (exit 0;
  \texttt{silent\_losses == 0}).
  \item Six deterministic MCP tools (I1): the registered tools in
  \pathtt{api/mcp/server.py} (\texttt{resolve\_account}, \texttt{get\_account},
  \texttt{validate\_balance\_sheet}, \texttt{consolidate\_trial\_balances},
  \texttt{get\_node\_fiber}, \texttt{list\_jurisdictions}); none invoke an LLM.
  \item Fiber / rollup (I3): \texttt{node\_fiber} and \texttt{rollup} in
  \pathtt{core/harness/ontology.py}; surfaced at REST
  \texttt{GET /accounts/\{id\}/fiber} and MCP \texttt{get\_node\_fiber}.
  \item $97.3\%$ deterministic resolution, $75$ entities, $68$ jurisdictions,
  $4$ escalations, and the tier distribution ($310$ Tier-1, $46$ Tier-2 over
  $366$ counted entries) quoted in Section~\ref{sec:inference-cost}:\\
  \texttt{python scripts/mass\_consolidation\_v2.py} $\rightarrow$
  \pathtt{research/experiments/consolidation_v2/results.json}
  (\texttt{tier\_distribution}, \texttt{deterministic\_coverage\_pct}).
  The $0.9\%$ / $99.1\%$ split in the same section is
  $4/441$ and $437/441$ from the round-trip artifact above.
  \item Figure~\ref{fig:fiber-query} (the $57$ local codes across $56$
  jurisdictions, and every code, name and local parent shown):\\
  \texttt{python}
  \pathtt{docs/papers/spokes/agentic-provenance/figures/gen_fig_fiber_query.py}\\
  regenerates the figure source directly from \texttt{node\_fiber} over the
  committed ontology, so the figure cannot drift from the data it depicts.
  Figure~\ref{fig:agent-flow} is structural; each of its labels is annotated
  in the figure source with the file and line it is taken from.
\end{itemize}

The design decision this paper reports is recorded as ADR-016
(\pathtt{docs/adr/016-lossless-translation-and-provenance.md}). The normative
specification is the canonical preprint, DOI 10.5281/zenodo.20738795.

\section*{Citing this work}

Cite Kontablo via its concept DOI, which resolves to the latest version:

\begin{quote}\small
Luciani, C. \emph{Kontablo: A Graph-Based Universal Accounting Ontology for the
M2M Agentic Economy.} Zenodo. \href{https://doi.org/10.5281/zenodo.20738795}{https://doi.org/10.5281/zenodo.20738795}
\end{quote}

\noindent This spoke paper is a community-facing companion; the concept DOI is
the citable specification of record. In BibTeX:

\begin{quote}\footnotesize
\begin{verbatim}
@misc{kontablo,
  author    = {Luciani, Christian},
  title     = {{Kontablo: A Graph-Based Universal Accounting
               Ontology for the M2M Agentic Economy}},
  year      = {2026},
  publisher = {Zenodo},
  doi       = {10.5281/zenodo.20738795},
  note      = {Concept DOI; resolves to the latest version},
  url       = {https://doi.org/10.5281/zenodo.20738795}
}
\end{verbatim}
\end{quote}

\noindent The reference implementation, all generating commands quoted above,
and the source of this paper are in the public repository
\href{https://github.com/ChristianLuciani/accounting-esperanto}{github.com/ChristianLuciani/accounting-esperanto}
(paper sources under \pathtt{docs/papers/spokes/agentic-provenance/}). This
spoke receives its own version DOI on publication; until then, cite the concept
DOI above.

\section*{License}

This paper --- its text, its figures, and the \LaTeX{} sources of both --- is
licensed \textbf{Creative Commons Attribution 4.0 International (CC BY 4.0)},
\href{https://creativecommons.org/licenses/by/4.0/}{creativecommons.org/licenses/by/4.0/},
consistent with the rest of the Kontablo paper corpus
(\pathtt{docs/papers/LICENSE}). You may share and adapt it, including
commercially, with attribution.

The license covers the paper, not the system it describes. The reference
implementation is licensed separately: the core under the Business Source
License 1.1, converting to Apache 2.0; the open-source ERP connectors under
Apache 2.0. Reusing this text does not grant any right to the implementation.

% =====================================================================
\bibliographystyle{plain}
\bibliography{references}
% =====================================================================

\end{document}
